% assessment - formative exercise student sheet (in-class, Fri 17 July 2026)
% IR227 Politics of Change: Governing Social Transformation
% ~30-minute in-class exercise mirroring the exam-question anatomy.
% In-class (seminar), ungraded. 
% The scenario is fictional.
\documentclass[nobib]{tufte-handout}
\usepackage{../shared/preamble/tufte-style}
\usepackage{../shared/macros}
\addbibresource{../literature/ir227-bib.bib}

\title{Formative Exercise: A Policy Window in Country F}
\author{IR227 Politics of Change: in-class formative (ungraded)}
\date{Friday 17 July 2026}

\begin{document}

\maketitle

\begin{abstract}
\noindent
This exercise builds the analytical skills the final examination rewards, and shows you how an exam question is structured before you sit the exam. The question below has the same anatomy as a final-exam question: a scenario, then parts asking you to apply, analyse, and evaluate. It is not graded. Bullet points are fine; full prose is not expected in thirty minutes. You will assess your own answer against the criteria overleaf, and we will debrief together.
\end{abstract}

\section{The scenario}

For more than a decade, Country F's plans to move its electricity system off imported gas and domestic coal went nowhere. The state power utility warned of blackout risk, coal-region legislators defended local jobs, and successive energy ministers judged the fight not worth having. Two reform bills died in committee between 2015 and 2022, and the issue stayed inside the energy ministry, attracting little national attention.

Over a fourteen-month period, the following developments occurred:

\begin{itemize}\itemsep=2pt
\item The national academy of engineering published analysis showing that gas-and-coal dependence had left Country F with some of the region's highest electricity prices and its greatest exposure to supply shocks.
\item A dispute with a gas supplier cut imports through a cold winter, causing weeks of rolling blackouts and a doubling of household energy bills. The crisis became the largest domestic news story of the year.
\item A general election returned a government that depended on a small green party for its majority, and the coalition agreement committed it to an energy transition.
\item A newly appointed energy-transition commissioner, a former national grid-operator chief who had become a public figure during the blackouts, was given a cross-ministry mandate and an explicit reform brief.
\item Energy-agency files, prepared years earlier and repeatedly shelved, contained a fully costed transition package modelled on systems already operating in over thirty countries: a binding renewables target, grid modernisation, a coal phase-out schedule, and an independent system operator.
\end{itemize}

Within the year, the government adopted the package, over the objections of the state utility and the coal-region legislators.

\section{The question}

\begin{enumerate}
\item[(a)] \textbf{Three streams:} Using the Multiple Streams Framework, identify the problem stream, the policy stream, and the political stream in this case. Which specific developments in the scenario constitute each stream?
\item[(b)] \textbf{Window and coupling:} Analyse how and why a policy window opened. Was it a predictable or an unpredictable window? What was the coupling, and who acted as the policy entrepreneur?
\item[(c)] \textbf{Evaluation:} Identify one thing MSF explains well about this case, and one aspect of the change the framework struggles to account for.
\end{enumerate}

Suggested timing: five minutes reading and planning, fifteen minutes writing, ten minutes for the self-check and paired comparison below.

\section{Self-check}

Score yourself honestly against each criterion, then compare answers with a neighbour before the plenary debrief.

\begin{description}\itemsep=3pt
\item[(a) Streams.] Did you tie \emph{each} stream to named developments in the scenario, rather than describing the streams in the abstract? Did anything in the scenario end up in the wrong stream?
\item[(b) Window.] Did you name the window and classify it (predictable or unpredictable)? Did you describe coupling as an \emph{action} that joined the streams, not just a coincidence? Did you identify the entrepreneur through what they \emph{invested}, not just who was prominent?
\item[(c) Evaluation.] Is your limitation a specific feature of the framework meeting this case (for instance, what MSF says about \emph{why this design} was selected, or whose resources counted), rather than generic scepticism?
\end{description}

\emph{The final exam asks for the same three moves at greater depth: about 60 minutes and 800 words per question, in full prose, answering two questions across the paper. The mid-term take-home essay uses the same analytical move on a case of your choosing; the one added expectation there is that you ground your analysis in the research literature and reference it, which the exam does not require.}

\end{document}
